% !TEX program = platex
\documentclass[a4j,12pt,openany]{jsreport}

% ===== パッケージ =====
\usepackage[dvipdfmx]{graphicx}
\usepackage[dvipdfmx]{xcolor}
\usepackage[dvipdfmx,bookmarksnumbered=true,colorlinks=true,
  linkcolor=black,urlcolor=black,citecolor=black,
  anchorcolor=black,filecolor=black]{hyperref}
\usepackage{pxjahyper}
\usepackage{listings}
\usepackage{jlisting}
\usepackage{geometry}
\usepackage{fancyhdr}
\usepackage{booktabs}
\usepackage{array}
\usepackage{enumitem}
\usepackage{amsmath}
\usepackage[skins,breakable]{tcolorbox}
\usepackage{url}
\usepackage{multirow}
\usepackage{tabularx}
\usepackage{longtable}
\usepackage{titlesec}

% ===== ページレイアウト =====
\geometry{top=20mm, bottom=20mm, left=20mm, right=20mm}

% ===== 表の可読性設定 =====
\renewcommand{\arraystretch}{1.35}
\setlength{\tabcolsep}{6pt}

% ===== 色定義 =====
\definecolor{bashdark}{HTML}{1E1E1E}
\definecolor{bashtext}{HTML}{D4D4D4}
\definecolor{psbg}{HTML}{012456}
\definecolor{pstext}{HTML}{EEEDF0}
\definecolor{notebg}{HTML}{EFF8E2}
\definecolor{noteframe}{HTML}{4CAF50}
\definecolor{warnbg}{HTML}{FFF8E1}
\definecolor{warnframe}{HTML}{FF8F00}
\definecolor{tipbg}{HTML}{E3F2FD}
\definecolor{tipframe}{HTML}{1565C0}
\definecolor{dangerbg}{HTML}{FFEBEE}
\definecolor{dangerframe}{HTML}{C62828}
\definecolor{pybg}{HTML}{F5F5FF}
\definecolor{pyframe}{HTML}{5C6BC0}

% ===== コードリスティング =====
\lstdefinestyle{bash}{
  backgroundcolor=\color{bashdark},
  basicstyle=\ttfamily\small\color{bashtext},
  keywordstyle=\color{cyan!80},
  commentstyle=\color{gray!70}\itshape,
  stringstyle=\color{yellow!70},
  showstringspaces=false,
  breaklines=true,
  frame=none,
  xleftmargin=1em,
  xrightmargin=1em,
  aboveskip=0.6em,
  belowskip=0.6em,
  language=bash,
}
\lstdefinestyle{python}{
  backgroundcolor=\color{pybg},
  basicstyle=\ttfamily\small,
  keywordstyle=\color{blue!70}\bfseries,
  commentstyle=\color{green!60!black}\itshape,
  stringstyle=\color{orange!70!red},
  showstringspaces=false,
  breaklines=true,
  frame=single,
  rulecolor=\color{pyframe},
  xleftmargin=1em,
  xrightmargin=1em,
  aboveskip=0.6em,
  belowskip=0.6em,
  language=python,
}
\lstdefinestyle{toml}{
  backgroundcolor=\color{pybg},
  basicstyle=\ttfamily\small,
  commentstyle=\color{green!60!black}\itshape,
  stringstyle=\color{orange!70!red},
  showstringspaces=false,
  breaklines=true,
  frame=single,
  rulecolor=\color{pyframe},
  xleftmargin=1em,
  xrightmargin=1em,
  aboveskip=0.6em,
  belowskip=0.6em,
}
\lstset{
  basicstyle=\ttfamily\small,
  breaklines=true,
}

% ===== tcolorbox 環境 =====
\newtcolorbox{wslbox}[1][Ubuntu/WSL ターミナル]{
  colback=bashdark,
  coltext=bashtext,
  coltitle=white,
  colbacktitle=black!60,
  fonttitle=\small\bfseries\ttfamily,
  title=#1,
  boxrule=0pt,
  sharp corners,
  left=0.5em, right=0.5em, top=0.3em, bottom=0.3em,
  titlerule=0pt,
  attach boxed title to top left={yshift=-0.1em},
  boxed title style={colback=black!60, sharp corners, boxrule=0pt},
  breakable,
}
\newtcolorbox{notebox}[1][メモ]{
  colback=notebg,
  colframe=noteframe,
  coltitle=white,
  colbacktitle=noteframe,
  fonttitle=\small\bfseries,
  title=#1,
  sharp corners,
  boxrule=0.5pt,
  breakable,
}
\newtcolorbox{warnbox}[1][注意]{
  colback=warnbg,
  colframe=warnframe,
  coltitle=white,
  colbacktitle=warnframe,
  fonttitle=\small\bfseries,
  title=#1,
  sharp corners,
  boxrule=0.5pt,
  breakable,
}
\newtcolorbox{tipbox}[1][ポイント]{
  colback=tipbg,
  colframe=tipframe,
  coltitle=white,
  colbacktitle=tipframe,
  fonttitle=\small\bfseries,
  title=#1,
  sharp corners,
  boxrule=0.5pt,
  breakable,
}
\newtcolorbox{dangerbox}[1][危険]{
  colback=dangerbg,
  colframe=dangerframe,
  coltitle=white,
  colbacktitle=dangerframe,
  fonttitle=\small\bfseries,
  title=#1,
  sharp corners,
  boxrule=0.5pt,
  breakable,
}

% ===== インラインコマンド =====
\newcommand{\cmd}[1]{\texttt{#1}}
\newcommand{\key}[1]{\fbox{\ttfamily\small #1}}
\newcommand{\win}{【Windows】}
\newcommand{\wsl}{【WSL】}

% ===== ヘッダ・フッタ =====
\pagestyle{fancy}
\fancyhf{}
\fancyhead[LE,RO]{\small\textit{B4課題 仮想環境セットアップガイド}}
\fancyhead[RE,LO]{\small\leftmark}
\fancyfoot[C]{\thepage}
\renewcommand{\headrulewidth}{0.4pt}

% ===== 章タイトルのスタイル =====
\titleformat{\chapter}[block]
  {\Large\bfseries}
  {\chaptertitlename\thechapter}
  {1em}{}
  [\vspace{0.5em}\hrule\vspace{0.5em}]

% ===== ドキュメント開始 =====
\begin{document}

% --- タイトルページ ---
\begin{titlepage}
  \vspace*{3cm}
  \begin{center}
    {\Huge\bfseries B4課題\\仮想環境セットアップガイド}\\[1.5em]
    {\Large uv を使った Python 環境構築と\\JupyterLab の起動}\\[3em]
    \hrule\vspace{1em}
    {\large
      前提：マニュアル第1弾（WSL 入門）完了済み\\[0.5em]
      2026 ver.
    }\\[2em]
    \hrule
  \end{center}
  \vfill
\end{titlepage}

% --- 目次 ---
\setcounter{tocdepth}{2}
\tableofcontents
\clearpage

% ===== 本文 =====
\chapter{はじめに}

このガイドは、\textbf{マニュアル第1弾（WSL 入門）の続編}です。
WSL（Ubuntu）のインストールとターミナルの基本操作は完了している前提で進めます。

\section{このガイドでやること}

\begin{itemize}
  \item 仮想環境とは何か、なぜ必要かを理解する
  \item パッケージマネージャ \textbf{uv} をインストールする
  \item \cmd{git clone} で課題リポジトリを取得する
  \item \cmd{uv sync} で必要なライブラリをまとめてインストールする
  \item VS Code でノートブックを開いて課題を実行する
\end{itemize}

\section{ゴール}

このガイドを読み終えると、VS Code で課題フォルダを開くだけで
作業を始められるようになります。

\begin{wslbox}[毎回の起動（最終目標）]
\begin{lstlisting}[style=bash]
$ code ~/NMR-lab-kadai/2026年度_B4課題引継ぎ
\end{lstlisting}
\end{wslbox}

\section{前提条件}

\begin{itemize}
  \item WSL（Ubuntu）がインストール済みであること
  \item WSL のターミナルを開けること
  \item インターネットに接続できること（初回セットアップ時のみ）
  \item VS Code に Jupyter 拡張機能がインストールされていること
\end{itemize}

\begin{notebox}
\textbf{Python のインストールは不要です。}\\
uv が Python 本体も含めて自動でダウンロード・管理します。
\end{notebox}

\chapter{仮想環境とは何か}

\section{なぜ仮想環境が必要か}

Python でプログラムを書くとき、\cmd{numpy} や \cmd{matplotlib} などの
外部ライブラリ（パッケージ）を使います。
これらには\textbf{バージョン}があり、プロジェクトによって必要なバージョンが異なります。

\subsection{バージョン衝突の問題}

次のようなシナリオを考えてみましょう。

\begin{center}
\begin{tabularx}{\linewidth}{lXX}
  \toprule
  & B4課題（NMR ラボ） & 別の研究プログラム \\
  \midrule
  numpy & 1.26 が必要 & 2.0 が必要 \\
  \bottomrule
\end{tabularx}
\end{center}

同じ Python 環境に両方をインストールしようとすると、
片方が上書きされてどちらかが動かなくなります。

\begin{center}
\begin{tabularx}{\linewidth}{lX}
  \toprule
  状況 & 結果 \\
  \midrule
  numpy 1.26 をインストール & 別の研究プログラムが動かない \\
  numpy 2.0 にアップグレード & B4課題が動かなくなる \\
  \bottomrule
\end{tabularx}
\end{center}

これが\textbf{依存関係の衝突}です。

\subsection{グローバル環境を汚さないために}

\cmd{pip install} を繰り返すと、システム全体の Python 環境に
パッケージがどんどん蓄積されます。

\begin{itemize}
  \item 使わなくなったパッケージが残り続ける
  \item バージョンアップで他のプログラムが突然壊れる
  \item 「先輩の PC では動いたのに自分の PC では動かない」が起きやすい
\end{itemize}

\section{仮想環境の仕組み}

\textbf{仮想環境（Virtual Environment）}はプロジェクトごとに
独立した Python の箱を作る仕組みです。

\begin{center}
\fbox{\parbox{0.85\textwidth}{
\textbf{仮想環境のイメージ}\\[0.5em]
システム全体\\
\quad└ Ubuntu のシステム Python（素の状態）\\
\quad\quad ├ B4課題用 仮想環境（\cmd{.venv}）\\
\quad\quad │ \quad numpy 1.26, bm3d 4.0, jupyterlab 4.0 ...\\
\quad\quad └ 別プロジェクト用 仮想環境（\cmd{.venv}）\\
\quad\quad \quad \quad numpy 2.0, tensorflow ...
}}
\end{center}

各仮想環境は\textbf{完全に独立}しており、一方を変更しても他方には影響しません。

\section{仮想環境の実体}

仮想環境はただのフォルダです。
プロジェクトフォルダの中に \cmd{.venv/} というフォルダが作られ、
その中にライブラリの実体が格納されます。

\begin{wslbox}[.venv の中身のイメージ]
\begin{lstlisting}[style=bash]
2026年度_B4課題引継ぎ/
  ├── pyproject.toml   # どのライブラリが必要か書いてある
  ├── uv.lock          # バージョンを固定するファイル
  ├── .venv/           # 仮想環境（uv sync で自動生成）
  │     ├── bin/
  │     └── lib/python3.11/site-packages/
  │           ├── numpy/
  │           ├── matplotlib/
  │           └── ...
  ├── No.1_FID/
  └── ...
\end{lstlisting}
\end{wslbox}

\begin{tipbox}
\cmd{.venv/} の中身は自動生成されるため、自分で触る必要はありません。
\cmd{uv sync} を実行するだけで作られます。
\end{tipbox}

\section{仮想環境を使うメリット}

\begin{description}
  \item[\textbf{依存関係の分離}]
    プロジェクトごとに必要なパッケージだけをインストールできる
  \item[\textbf{再現性}]
    全員が同じバージョンのライブラリを使えるので
    「自分の環境では動くのに\ldots\ldots」という問題が起きない
  \item[\textbf{削除が簡単}]
    不要になったら \cmd{.venv/} フォルダを削除するだけ
  \item[\textbf{システムを汚さない}]
    Ubuntu のシステム Python は変更しないので OS が安定する
\end{description}

\chapter{uv の仕組みを理解する}

\section{uv の概要}

\textbf{uv}（読み方：ユーブイ）は Astral 社が開発した
Python パッケージマネージャです。
公式サイト：\url{https://docs.astral.sh/uv/}

従来の \cmd{pip}（パッケージインストーラー）と
\cmd{venv}（仮想環境作成ツール）の両方の機能を一つに統合しており、
以下の特徴があります。

\begin{description}
  \item[\textbf{高速}]
    Rust で書かれており、\cmd{pip} より 10〜100 倍速い。
    多数のパッケージをインストールしても数秒で完了する
  \item[\textbf{再現性が高い}]
    \cmd{uv.lock} ファイルで全員がまったく同じバージョンを使える
  \item[\textbf{Python も自動管理}]
    Python 本体のインストールも自動で行う。
    事前に Python をインストールしておく必要がない
  \item[\textbf{シンプル}]
    \cmd{uv sync} 一発で仮想環境の作成からパッケージのインストールまで完了する
\end{description}

\section{pip + venv との比較}

従来の方法（\cmd{pip} + \cmd{venv}）と \cmd{uv} を比べると次のとおりです。

\begin{center}
\begin{tabularx}{\linewidth}{lXX}
  \toprule
  操作 & pip + venv & uv \\
  \midrule
  Python のインストール & 別途必要 & 自動 \\
  仮想環境を作る & \cmd{python3 -m venv .venv} & 自動（\cmd{uv sync} で） \\
  仮想環境を有効化 & \cmd{source .venv/bin/activate} & 不要 \\
  パッケージをインストール & \cmd{pip install -r req.txt} & \cmd{uv sync} \\
  スクリプトを実行 & \cmd{python script.py} & \cmd{uv run python script.py} \\
  パッケージを追加 & \cmd{pip install xxx} & \cmd{uv add xxx} \\
  新規プロジェクト作成 & 手動で \cmd{pyproject.toml} 作成 & \cmd{uv init} \\
  速度 & 普通 & 10〜100 倍速い \\
  \bottomrule
\end{tabularx}
\end{center}

\begin{tipbox}[activate が不要な理由]
uv では \cmd{source .venv/bin/activate} による仮想環境の有効化が不要です。
\cmd{uv run} をつけるだけで、自動的に \cmd{.venv/} 内の Python と
ライブラリを使って実行してくれます。
「有効化し忘れてシステム Python で実行してしまった」という
よくあるミスがなくなります。
\end{tipbox}

% -------------------------------------------------
\section{pyproject.toml とは}
% -------------------------------------------------

\cmd{pyproject.toml} はプロジェクトの設定ファイルです。
「このプロジェクトには何が必要か」をまとめて書く場所です。

\subsection{ファイルの構造}

B4課題フォルダの \cmd{pyproject.toml} を見てみましょう。

\begin{lstlisting}[style=toml, caption={pyproject.toml の構造}]
[project]
name = "nmrlab-b4-2026"      # プロジェクトの名前
version = "2026.1"            # バージョン
requires-python = ">=3.11"   # Python 3.11 以上が必要

dependencies = [             # 必要なライブラリの一覧
    "numpy>=1.26",           # numpy バージョン 1.26 以上
    "scipy>=1.12",
    "matplotlib>=3.8",
    "scikit-image>=0.22",
    "PyWavelets>=1.5",
    "bm3d>=4.0",
    "openpyxl>=3.1",
    "jupyterlab>=4.0",
    "ipywidgets>=8.0",
    "japanize-matplotlib>=1.1",
]

[tool.uv]
dev-dependencies = [         # 開発時だけ使うツール（通常不要）
    "pytest>=8.0",
]
\end{lstlisting}

\subsection{バージョン指定の書き方}

\begin{center}
\begin{tabularx}{\linewidth}{lX}
  \toprule
  書き方 & 意味 \\
  \midrule
  \cmd{"numpy>=1.26"} & 1.26 以上ならどのバージョンでも OK \\
  \cmd{"numpy==1.26.4"} & 1.26.4 ぴったり（固定） \\
  \cmd{"numpy>=1.26,<2.0"} & 1.26 以上かつ 2.0 未満 \\
  \cmd{"numpy"} & バージョン指定なし（最新） \\
  \bottomrule
\end{tabularx}
\end{center}

\begin{notebox}[pyproject.toml は「希望」を書く場所]
\cmd{pyproject.toml} の \cmd{dependencies} は
「このバージョン以上なら動く」という\textbf{条件}を書くところです。
実際にインストールする\textbf{正確なバージョン}は
\cmd{uv.lock} に記録されます。
\end{notebox}

\subsection{パッケージを追加するには}

\cmd{pyproject.toml} を直接編集してもよいですが、
\cmd{uv add} コマンドを使うと自動で追記されます。

\begin{wslbox}[パッケージの追加]
\begin{lstlisting}[style=bash]
$ uv add pandas
\end{lstlisting}
\end{wslbox}

実行すると \cmd{pyproject.toml} の \cmd{dependencies} に
\cmd{"pandas>=X.X.X"} が自動追記されます。

% -------------------------------------------------
\section{uv.lock とは}
% -------------------------------------------------

\cmd{uv.lock} は uv が自動生成する\textbf{ロックファイル}です。
\textbf{絶対に手動で編集してはいけません。}

\subsection{なぜ必要か}

\cmd{pyproject.toml} には「numpy 1.26 以上」と書きましたが、
これだけでは「今インストールしたら 1.26.4 なのか 1.27.0 なのか」が
わかりません。

時間が経つと最新バージョンが変わるため、
同じ \cmd{pyproject.toml} を使っても
インストールのタイミングによって入るバージョンが変わります。

\begin{center}
\fbox{\parbox{0.85\textwidth}{
\textbf{ロックファイルがないと起きる問題}\\[0.5em]
自分（3月にインストール）：numpy 1.26.4\\
後輩（9月にインストール）：numpy 1.27.2\\
→ 同じコードが異なる環境で動く可能性がある
}}
\end{center}

\subsection{uv.lock が解決すること}

\cmd{uv.lock} には\textbf{全パッケージの正確なバージョンとハッシュ値}が
記録されます。

\begin{wslbox}[uv.lock の中身（抜粋）]
\begin{lstlisting}[style=bash]
version = 1
requires-python = ">=3.11"

[[package]]
name = "numpy"
version = "1.26.4"    # 正確なバージョンが固定されている
source = { registry = "..." }
sdist = { url = "...", hash = "sha256:abc123..." }
\end{lstlisting}
\end{wslbox}

\cmd{uv sync} は \cmd{uv.lock} を読んで\textbf{全員が同じバージョン}を
インストールします。誰がいつ実行しても同じ環境になります。

\subsection{uv.lock の管理ルール}

\begin{itemize}
  \item \textbf{手動で編集しない}（uv が自動管理する）
  \item \textbf{共有フォルダやリポジトリに含める}（全員で共有する）
  \item \cmd{uv add}/\cmd{uv remove} を実行すると自動更新される
  \item \cmd{uv sync} を実行しても \cmd{uv.lock} は変化しない
    （\cmd{uv.lock} に従ってインストールするだけ）
\end{itemize}

\begin{tipbox}[pyproject.toml と uv.lock の関係まとめ]
\begin{tabularx}{\linewidth}{lXX}
  \toprule
  & \cmd{pyproject.toml} & \cmd{uv.lock} \\
  \midrule
  書いてあること & 「これ以上のバージョンが必要」という条件 & 実際にインストールする正確なバージョン \\
  誰が書く & 人間（uv add でも可） & uv が自動生成 \\
  手動編集 & OK & してはいけない \\
  共有する & する & する \\
  \bottomrule
\end{tabularx}
\end{tipbox}

% -------------------------------------------------
\section{uv sync の動作}
% -------------------------------------------------

\cmd{uv sync} は以下のことを\textbf{すべて自動で}行います。

\begin{enumerate}
  \item \cmd{pyproject.toml} を読んで必要な Python バージョンを確認
  \item 必要な Python がなければ自動ダウンロード・インストール
  \item \cmd{.venv/} フォルダを作成（すでにあれば再利用）
  \item \cmd{uv.lock} を読んで記録されたバージョンのパッケージを
    \cmd{.venv/} にインストール
  \item 余分なパッケージ（\cmd{uv.lock} にないもの）があれば削除
\end{enumerate}

\begin{wslbox}[uv sync の出力例（初回）]
\begin{lstlisting}[style=bash]
$ uv sync
Using CPython 3.11.x        <- Python を自動選択
Creating virtual environment at: .venv  <- .venv を作成
Resolved 47 packages in 0.5ms
Installed 47 packages in 8.2s  <- ここが pip より圧倒的に速い
 + bm3d==4.0.2
 + numpy==1.26.4
 + jupyterlab==4.3.6
 ...
\end{lstlisting}
\end{wslbox}

\begin{wslbox}[uv sync の出力例（2回目以降・差分なし）]
\begin{lstlisting}[style=bash]
$ uv sync
Resolved 47 packages in 0.5ms
Audited 47 packages in 1ms   <- 確認のみ、インストールなし
\end{lstlisting}
\end{wslbox}

\begin{notebox}[uv sync は「冪等」]
何度実行しても結果は同じです。
すでに正しい状態なら何もしません。
環境がおかしいと感じたらとりあえず \cmd{uv sync} を試してください。
\end{notebox}

\chapter{uv のインストール}

\section{インストールコマンド}

WSL のターミナルを開いて以下を実行します。
\textbf{これは最初の1回だけ}でよいです。

\begin{wslbox}[uv のインストール]
\begin{lstlisting}[style=bash]
$ curl -LsSf https://astral.sh/uv/install.sh | sh
\end{lstlisting}
\end{wslbox}

実行すると以下のような出力が出ます。

\begin{wslbox}[インストール時の出力例]
\begin{lstlisting}[style=bash]
downloading uv 0.6.x x86_64-unknown-linux-musl
installing to /home/username/.local/bin
  uv
  uvx
everything's installed!

To add $HOME/.local/bin to your PATH, either
restart your shell or run:

    source $HOME/.bashrc
\end{lstlisting}
\end{wslbox}

\begin{notebox}[\cmd{curl ... | sh} について]
インターネットからスクリプトをダウンロードして直接実行するコマンドです。
Astral 社の公式インストーラーなので安全です。
\end{notebox}

\section{PATH の反映}

インストール直後はコマンドがまだ認識されません。
\textbf{ターミナルを一度閉じて開き直してください。}

あるいは、ターミナルを開き直さずに即座に反映させる場合は：

\begin{wslbox}[PATH を今すぐ反映する]
\begin{lstlisting}[style=bash]
$ source ~/.bashrc
\end{lstlisting}
\end{wslbox}

\section{インストールの確認}

以下のコマンドでバージョンを確認します。

\begin{wslbox}[バージョン確認]
\begin{lstlisting}[style=bash]
$ uv --version
uv 0.6.x (linux x86_64 ...)
\end{lstlisting}
\end{wslbox}

このように表示されれば OK です。

\begin{wslbox}[失敗した場合の出力]
\begin{lstlisting}[style=bash]
$ uv --version
-bash: uv: command not found
\end{lstlisting}
\end{wslbox}

\cmd{command not found} と表示された場合は、
ターミナルの再起動か \cmd{source ~/.bashrc} を試してください。
それでも解決しない場合は第~\ref{chap:errors}章を参照してください。

\section{uv のアップデート（将来の参考用）}

uv 自体のアップデートは以下のコマンドで行えます。
セットアップ時に実行する必要はありません。

\begin{wslbox}[uv のアップデート]
\begin{lstlisting}[style=bash]
$ uv self update
\end{lstlisting}
\end{wslbox}

\chapter{フォルダ構成と配置場所}

\section{配置先のディレクトリ}

\cmd{git clone} を実行すると、ホームディレクトリ直下に
リポジトリフォルダが作成されます。
課題フォルダはリポジトリ直下の \cmd{2026年度\_B4課題引継ぎ/} にあります。

\begin{wslbox}[フォルダ構成]
\begin{lstlisting}[style=bash]
/home/username/              <- ホームディレクトリ（~ と同じ）
└── NMR-lab-kadai/           <- git clone で作成される
    ├── 2026年度_B4課題引継ぎ/  <- VS Code で開くフォルダ
    │   ├── .vscode/
    │   │   └── settings.json  <- カーネル自動設定
    │   ├── pyproject.toml
    │   ├── uv.lock
    │   ├── No.1_FID/
    │   │   ├── makesignal.ipynb
    │   │   └── heikatsu.ipynb
    │   ├── No.2_FFT1D/
    │   │   ├── FFT1D.ipynb
    │   │   ├── kukei.ipynb
    │   │   └── sft.ipynb
    │   ├── No.3_FFT2D/
    │   │   └── FFT2D.ipynb
    │   ├── No.4_DFT_exercise/
    │   │   └── DFT_template.ipynb
    │   ├── No.5_CompressedSensing/
    │   │   ├── cs_tv.ipynb
    │   │   └── cs_wavelet.ipynb
    │   ├── No.6_Denoising/
    │   │   ├── add_noise.ipynb
    │   │   └── denoise_bm3d.ipynb
    │   ├── No.7_Metrics/
    │   │   └── psnr_ssim.ipynb
    │   └── No.8_CT/
    │       └── ct_pipeline.ipynb
    └── uv-manual/             <- セットアップマニュアル
\end{lstlisting}
\end{wslbox}

\begin{notebox}[username について]
\cmd{username} の部分は自分の WSL ユーザー名です。
次のコマンドで確認できます。

\begin{wslbox}[ユーザー名の確認]
\begin{lstlisting}[style=bash]
$ whoami
\end{lstlisting}
\end{wslbox}
\end{notebox}

\section{重要なファイルの説明}

\begin{center}
\begin{tabularx}{\linewidth}{lX}
  \toprule
  ファイル & 役割 \\
  \midrule
  \cmd{pyproject.toml} &
    必要なライブラリとバージョン条件が書かれた設定ファイル。
    \cmd{uv sync} はこれを読んでインストール内容を決める \\
  \cmd{uv.lock} &
    実際にインストールするバージョンを厳密に記録したファイル。
    全員が同じ環境を再現するために使う。
    \textbf{編集しないこと} \\
  \cmd{.venv/} &
    \cmd{uv sync} を実行すると自動生成される仮想環境フォルダ。
    git には含まれないので各自が \cmd{uv sync} で作成する \\
  \cmd{.vscode/settings.json} &
    VS Code が自動で \cmd{.venv} の Python を選ぶための設定ファイル。
    カーネル選択の手間を省く \\
  \cmd{*.ipynb} &
    JupyterLab で開く課題ノートブック \\
  \cmd{*.py} &
    対応する Python スクリプト（ノートブックと同じ処理） \\
  \bottomrule
\end{tabularx}
\end{center}

\begin{notebox}[.venv/ は git 管理外]
\cmd{.venv/} フォルダはサイズが数百 MB あり、
\cmd{uv sync} で誰でも再生成できるため git には含まれていません。
\cmd{git clone} 後に各自で \cmd{uv sync} を実行してください。
\end{notebox}

\chapter{課題リポジトリを取得する}

課題ファイルは GitHub 上で管理されています。
\cmd{git clone} コマンドで WSL のホームディレクトリに取得します。
\textbf{この操作は最初の1回だけ}行います。

% =================================================
\section{git clone で取得する}
% =================================================

WSL ターミナルで以下を実行します。

\begin{wslbox}[リポジトリのクローン]
\begin{lstlisting}[style=bash]
$ cd ~
$ git clone \
  https://github.com/shiru2/NMR-lab-kadai.git
\end{lstlisting}
\end{wslbox}

正常に完了すると \cmd{NMR-lab-kadai} フォルダが作成されます。

\begin{wslbox}[取得の確認]
\begin{lstlisting}[style=bash]
$ ls \
  ~/NMR-lab-kadai/2026年度_B4課題引継ぎ/
No.1_FID  No.2_FFT1D  No.3_FFT2D
No.4_DFT_exercise  No.5_CompressedSensing
No.6_Denoising  No.7_Metrics  No.8_CT
README.md  pyproject.toml  uv.lock
\end{lstlisting}
\end{wslbox}

\cmd{pyproject.toml} と \cmd{uv.lock} が表示されれば正しく取得できています。

\begin{notebox}[git が使えない場合]
大学のネットワーク環境によって \cmd{git clone} が
ブロックされる場合があります。
その場合は GitHub のページから zip をダウンロードして展開してください。\\[0.5em]
\texttt{Code} $\rightarrow$ \texttt{Download ZIP} を選び、
展開したフォルダを\\
\cmd{{\textbackslash}{\textbackslash}wsl.localhost{\textbackslash}Ubuntu{\textbackslash}home{\textbackslash}username}\\
に配置してください（\cmd{username} は自分の WSL ユーザー名）。
\end{notebox}

% =================================================
\section{2回目以降：最新版に更新する}
% =================================================

課題ファイルが更新された場合は \cmd{git pull} で最新版を取得できます。

\begin{wslbox}[最新版に更新]
\begin{lstlisting}[style=bash]
$ cd ~/NMR-lab-kadai
$ git pull
\end{lstlisting}
\end{wslbox}

\chapter{仮想環境のセットアップ}

\section{フォルダに移動する}

まず課題フォルダに移動します。

\begin{wslbox}[課題フォルダへ移動]
\begin{lstlisting}[style=bash]
$ cd \
  ~/NMR-lab-kadai/2026_B4_assignment
\end{lstlisting}
\end{wslbox}

移動できたか確認します。

\begin{wslbox}[現在地の確認]
\begin{lstlisting}[style=bash]
$ pwd
/home/username/NMR-lab-kadai/2026_B4_assignment
\end{lstlisting}
\end{wslbox}

このように表示されれば OK です。

\section{uv sync を実行する}

\begin{wslbox}[仮想環境のセットアップ]
\begin{lstlisting}[style=bash]
$ uv sync
\end{lstlisting}
\end{wslbox}

初回実行時は Python のダウンロードと全ライブラリのインストールが
自動で行われます。回線速度によっては数分かかる場合があります。

\subsection{実行中の出力例}

\begin{wslbox}[uv sync の出力例]
\begin{lstlisting}[style=bash]
Using CPython 3.11.x
Creating virtual environment at: .venv
Resolved 47 packages in 0.5ms
Installed 47 packages in 8.2s
 + bm3d==4.0.2
 + ipywidgets==8.1.5
 + japanize-matplotlib==1.1.3
 + jupyterlab==4.3.6
 + matplotlib==3.9.4
 + numpy==1.26.4
 + PyWavelets==1.5.0
 + scikit-image==0.22.0
 + scipy==1.13.1
 ...
\end{lstlisting}
\end{wslbox}

このような出力が出れば完了です。

\subsection{2回目以降の実行}

すでに環境が揃っている場合は、\cmd{uv sync} を再実行しても
差分だけが更新されます。不足しているパッケージがあれば追加されます。

\begin{wslbox}[2回目以降の出力例]
\begin{lstlisting}[style=bash]
$ uv sync
Resolved 47 packages in 0.5ms
Audited 47 packages in 1ms
\end{lstlisting}
\end{wslbox}

\section{セットアップの確認}

正しくインストールされたか確認します。

\begin{wslbox}[インストール済みパッケージの確認]
\begin{lstlisting}[style=bash]
$ uv pip list
Package                  Version
------------------------ -------
bm3d                     4.0.2
japanize-matplotlib      1.1.3
jupyterlab               4.3.6
matplotlib               3.9.4
numpy                    1.26.4
...
\end{lstlisting}
\end{wslbox}

\cmd{numpy}、\cmd{jupyterlab}、\cmd{bm3d} などが表示されれば
正しくインストールされています。

\begin{tipbox}[uv.lock の役割]
\cmd{uv.lock} ファイルがあることで、
先輩・後輩・教員の全員がまったく同じバージョンのライブラリを使えます。
「自分の環境では動いたのに先輩の環境では動かなかった」という問題を防ぐための仕組みです。
\end{tipbox}

\chapter{ノートブックの実行方法}

% =================================================
\section{VS Code でノートブックを開く（推奨）}
% =================================================

\subsection{事前準備：Jupyter 拡張機能のインストール}

VS Code を初めて使う場合のみ必要です。

\begin{enumerate}
  \item VS Code を開く
  \item 左側の拡張機能アイコン（\key{Ctrl}+\key{Shift}+\key{X}）をクリック
  \item 検索欄に \texttt{Jupyter} と入力
  \item \texttt{Jupyter}（Microsoft 製）をインストール
\end{enumerate}

\subsection{フォルダを開いてノートブックを実行する}

\begin{enumerate}
  \item VS Code で \cmd{File} $\rightarrow$ \cmd{Open Folder} を選ぶ
  \item \cmd{2026年度\_B4課題引継ぎ} フォルダを開く
  \item 左側のファイルツリーから \cmd{.ipynb} ファイルをクリック
  \item 初回はカーネル選択が求められるので
    \cmd{.venv/bin/python} を選ぶ
\end{enumerate}

\begin{notebox}[カーネルの選び方]
ノートブックを開いたとき、右上に
「カーネルを選択」または Python のバージョンが表示されます。
クリックして一覧から \cmd{.venv/bin/python}（または
\texttt{Python 3.11.x ('.venv')}）を選んでください。\\[0.5em]
一覧に出ない場合は先に \cmd{uv sync} を実行してから
VS Code を開き直してください。
\end{notebox}

\subsection{セルの実行}

\begin{center}
\begin{tabularx}{\linewidth}{lX}
  \toprule
  操作 & 方法 \\
  \midrule
  セルを実行する & \key{Shift} + \key{Enter} \\
  セルを実行（その場にとどまる） & \key{Ctrl} + \key{Enter} \\
  全セルを一括実行 & ノートブック上部の \texttt{Run All} ボタン \\
  カーネルを再起動 & \texttt{Restart} ボタンまたは \key{F1} で検索 \\
  \bottomrule
\end{tabularx}
\end{center}

% =================================================
\section{課題ノートブックの一覧}
% =================================================

\begin{center}
\begin{tabularx}{\linewidth}{ll>{\raggedright\arraybackslash}X}
  \toprule
  課題 & ファイル & 内容 \\
  \midrule
  No.1 & \cmd{No.1\_FID/makesignal.ipynb} & FID 信号生成とノイズ付加 \\
  No.1 & \cmd{No.1\_FID/heikatsu.ipynb} & 移動平均による平滑化 \\
  No.2 & \cmd{No.2\_FFT1D/FFT1D.ipynb} & 1次元 FFT \\
  No.2 & \cmd{No.2\_FFT1D/kukei.ipynb} & 矩形パルスの FFT \\
  No.2 & \cmd{No.2\_FFT1D/sft.ipynb} & スライディング FFT \\
  No.3 & \cmd{No.3\_FFT2D/FFT2D.ipynb} & 2次元 FFT（MRI k空間） \\
  No.4 & \cmd{No.4\_DFT\_exercise/DFT\_template.ipynb} & DFT 手動実装（演習） \\
  No.5 & \cmd{No.5\_CompressedSensing/cs\_tv.ipynb} & 圧縮センシング（TV 法） \\
  No.5 & \cmd{No.5\_CompressedSensing/cs\_wavelet.ipynb} & 圧縮センシング（Wavelet） \\
  No.6 & \cmd{No.6\_Denoising/add\_noise.ipynb} & ノイズ付加 \\
  No.6 & \cmd{No.6\_Denoising/denoise\_bm3d.ipynb} & BM3D デノイジング \\
  No.7 & \cmd{No.7\_Metrics/psnr\_ssim.ipynb} & PSNR・SSIM 計算 \\
  No.8 & \cmd{No.8\_CT/ct\_pipeline.ipynb} & CT フィルタード逆投影 \\
  \bottomrule
\end{tabularx}
\end{center}

\begin{notebox}[No.4 は演習形式]
\cmd{DFT\_template.ipynb} は空白セルに自分でコードを書く演習です。
解答例は \cmd{DFT\_answer.ipynb} にあります。
まず自分で考えてから見るようにしてください。
\end{notebox}

\begin{notebox}[No.6 の WNNM はオプション]
\cmd{denoise\_wnnm.ipynb} は発展課題です。
256$\times$256 画像で数分かかるため、
まず \cmd{denoise\_bm3d.ipynb} を完了させてから取り組んでください。
\end{notebox}

% =================================================
\section{JupyterLab を使う場合（任意）}
% =================================================

ブラウザ上でノートブックを操作したい場合は JupyterLab を使えます。

\begin{wslbox}[JupyterLab の起動]
\begin{lstlisting}[style=bash]
$ cd ~/NMR-lab-kadai/2026年度_B4課題引継ぎ
$ uv run jupyter lab
\end{lstlisting}
\end{wslbox}

\begin{warnbox}[\cmd{uv run} を必ずつける]
\cmd{jupyter lab} だけで起動すると仮想環境の外から起動してしまい、
ライブラリが見つかりません。
必ず \cmd{uv run jupyter lab} と打ってください。
\end{warnbox}

起動するとターミナルに URL が表示されます。
ブラウザで \cmd{http://localhost:8888} を開いてください。
（WSL では自動でブラウザが開かない場合があります。
その場合は URL 全体をコピーして貼り付けてください。）

終了するときはターミナルで \key{Ctrl} + \key{C} を押します。

\begin{warnbox}[ブラウザのタブを閉じるだけでは終了しない]
タブを閉じてもサーバーは動き続けます。
必ずターミナルで \key{Ctrl} + \key{C} を押して終了してください。
\end{warnbox}

\chapter{よくあるエラーと対処}
\label{chap:errors}

\section{\cmd{uv: command not found}}

\textbf{原因}：uv がインストールされていないか、PATH が通っていません。

\textbf{対処1}：ターミナルを閉じて開き直す。

\textbf{対処2}：PATH を今すぐ反映させる。

\begin{wslbox}
\begin{lstlisting}[style=bash]
$ source ~/.bashrc
\end{lstlisting}
\end{wslbox}

\textbf{対処3}：それでも解決しない場合は PATH を手動で追加する。

\begin{wslbox}[PATH を手動で追加]
\begin{lstlisting}[style=bash]
$ echo 'export PATH="$HOME/.local/bin:$PATH"' \
  >> ~/.bashrc
$ source ~/.bashrc
\end{lstlisting}
\end{wslbox}

最後に確認：

\begin{wslbox}[確認]
\begin{lstlisting}[style=bash]
$ uv --version
uv 0.6.x (linux x86_64 ...)
\end{lstlisting}
\end{wslbox}

\section{\cmd{No pyproject.toml found}}

\textbf{原因}：\cmd{pyproject.toml} のあるフォルダにいません。

\textbf{対処}：正しいフォルダに移動してから実行する。

\begin{wslbox}[正しいフォルダへ移動して実行]
\begin{lstlisting}[style=bash]
$ cd ~/NMR-lab-kadai/2026_B4_assignment
$ uv sync
\end{lstlisting}
\end{wslbox}

フォルダに \cmd{pyproject.toml} があるか確認：

\begin{wslbox}[pyproject.toml の確認]
\begin{lstlisting}[style=bash]
$ ls pyproject.toml
pyproject.toml
\end{lstlisting}
\end{wslbox}

\cmd{pyproject.toml} と表示されれば正しいフォルダにいます。

\section{JupyterLab でカーネルが見つからない}

\textbf{原因}：\cmd{uv run} をつけずに JupyterLab を起動しています。

\textbf{対処}：ターミナルで \key{Ctrl} + \key{C} を押して一度終了し、
正しいコマンドで起動し直す。

\begin{wslbox}[正しい起動コマンド]
\begin{lstlisting}[style=bash]
$ cd ~/NMR-lab-kadai/2026_B4_assignment
$ uv run jupyter lab
\end{lstlisting}
\end{wslbox}

\section{\cmd{ModuleNotFoundError} が出る}

\textbf{原因1}：\cmd{uv sync} が完了していない。

\textbf{対処}：\cmd{uv sync} を再実行する。

\begin{wslbox}
\begin{lstlisting}[style=bash]
$ cd ~/NMR-lab-kadai/2026_B4_assignment
$ uv sync
\end{lstlisting}
\end{wslbox}

\textbf{原因2}：\cmd{uv run} なしで JupyterLab を起動している。
$\rightarrow$ 前節の対処と同じ手順で起動し直す。

\section{\cmd{uv sync} がネットワークエラーで止まる}

\textbf{原因}：インターネット接続が不安定、または
学内ネットワークのプロキシが干渉しています。

\textbf{対処}：数分待ってから再実行する。
\cmd{uv sync} は途中まで進んだ場合、次回実行時に続きから再開します。

\begin{wslbox}[再実行]
\begin{lstlisting}[style=bash]
$ uv sync
\end{lstlisting}
\end{wslbox}

繰り返し発生する場合は教員・先輩に相談してください。

\section{VS Code で「カーネルを起動できませんでした」と表示される}

\textbf{症状}：ノートブックを開くと
「Python 環境 'nmrlab-b4-2026 (3.12.x)' を使用できなくなったため、
カーネルを起動できませんでした」と表示される。

\textbf{原因}：VS Code が以前選択したカーネルをキャッシュしているため、
\cmd{.venv/} を作り直したあとに古い参照が残ります。
\cmd{pyvenv.cfg} を手動で書き換えても \cmd{uv sync} で元に戻るので、
この方法では解決しません。

\textbf{対処}：\cmd{uv sync} で \cmd{.venv/} を作り直したあと、
VS Code でカーネルを選び直す。

\begin{enumerate}
  \item ターミナルで \cmd{.venv/} を削除して再構築する
\end{enumerate}

\begin{wslbox}[.venv を削除して再構築]
\begin{lstlisting}[style=bash]
$ cd ~/NMR-lab-kadai/2026_B4_assignment
$ rm -rf .venv
$ uv sync
\end{lstlisting}
\end{wslbox}

\begin{enumerate}
  \setcounter{enumi}{1}
  \item VS Code でノートブックを開き、右上の
    「カーネルを選択」をクリック
  \item 「別のカーネルを選択」$\rightarrow$「Python 環境」を選ぶ
  \item 一覧から \cmd{.venv/bin/python}（または
    \texttt{nmrlab-b4-2026 (3.12.x)}）を選択する
\end{enumerate}

\begin{notebox}[一覧に出ない場合]
「Python 環境の一覧を更新」をクリックするか、
VS Code を開き直してから再度選択してください。
\end{notebox}

\section{Jupyter で「Python 3.12.x が使用できなくなった」と表示される}

\textbf{原因}：フォルダ内に古い \cmd{.venv/} が残っており、
そこに保存されていた Python のパスが現在の環境に存在しません。
（配布ファイルを古いバージョンから上書き展開した場合に起こります。）

\textbf{対処}：\cmd{.venv/} を削除してから \cmd{uv sync} を再実行する。

\begin{wslbox}[.venv を削除して再構築]
\begin{lstlisting}[style=bash]
$ cd ~/NMR-lab-kadai/2026_B4_assignment
$ rm -rf .venv
$ uv sync
\end{lstlisting}
\end{wslbox}

\cmd{uv sync} が完了したら VS Code でフォルダを開き直し、
カーネルを \cmd{.venv/bin/python} に選び直してください。

\section{\cmd{uv sync} 後に \cmd{Permission denied} が出る}

\textbf{原因}：フォルダを Windows の共有ドライブから直接コピーした場合、
ファイルの実行権限が失われることがあります。

\textbf{対処1}：\cmd{.venv/} を削除して再構築する（最も確実）。

\begin{wslbox}[.venv を削除して再構築]
\begin{lstlisting}[style=bash]
$ cd ~/NMR-lab-kadai/2026_B4_assignment
$ rm -rf .venv
$ uv sync
\end{lstlisting}
\end{wslbox}

\textbf{対処2}：フォルダ全体の権限を修正する。

\begin{wslbox}[権限の変更]
\begin{lstlisting}[style=bash]
$ chmod -R u+rw \
  ~/NMR-lab-kadai/2026_B4_assignment
\end{lstlisting}
\end{wslbox}

\begin{notebox}[git clone で取得した場合は発生しにくい]
\cmd{git clone} で正しく取得した場合、パーミッション問題は
ほとんど起きません。問題が続く場合は \cmd{.venv/} を削除して
\cmd{uv sync} を再実行してください（第6章参照）。
\end{notebox}

\section{ブラウザが自動で開かない}

\textbf{原因}：WSL から Windows のブラウザを自動起動できない場合があります。

\textbf{対処}：ターミナルに表示された URL を手動でコピーしてブラウザに貼り付ける。

\begin{wslbox}[URL の場所]
\begin{lstlisting}[style=bash]
[I ServerApp] Jupyter Server is running at:
[I ServerApp] http://localhost:8888/lab
  ?token=abc123def456...  <- これを全部コピーする
\end{lstlisting}
\end{wslbox}

\chapter{補足知識}

\section{\cmd{.venv/} フォルダの扱い}

\cmd{.venv/} は各自の PC で \cmd{uv sync} を実行すれば再生成できます。
以下のことを覚えておいてください。

\begin{itemize}
  \item 共有フォルダに含める必要はない
  \item メールや zip でやり取りする必要はない
  \item 誤って削除しても \cmd{uv sync} で復元できる
  \item リポジトリの \cmd{.gitignore} で除外済み
\end{itemize}

\section{環境を作り直す}

環境が壊れたと感じたときは \cmd{.venv/} を削除して作り直せます。

\begin{wslbox}[環境の再作成]
\begin{lstlisting}[style=bash]
$ cd ~/NMR-lab-kadai/2026_B4_assignment
$ rm -rf .venv
$ uv sync
\end{lstlisting}
\end{wslbox}

\section{新しいパッケージを追加する}

演習の範囲を超えて新しいライブラリを試したい場合：

\begin{wslbox}[パッケージの追加]
\begin{lstlisting}[style=bash]
$ uv add パッケージ名
\end{lstlisting}
\end{wslbox}

例えば \cmd{pandas} を追加する場合：

\begin{wslbox}[pandas の追加例]
\begin{lstlisting}[style=bash]
$ uv add pandas
Resolved 49 packages in 0.3ms
Installed 2 packages in 1.2s
 + pandas==2.2.3
 + python-dateutil==2.9.0
\end{lstlisting}
\end{wslbox}

\cmd{pyproject.toml} と \cmd{uv.lock} が自動的に更新されます。

\section{スクリプトを直接実行する}

ノートブックを使わず Python スクリプト（\cmd{.py} ファイル）を
直接実行したい場合も \cmd{uv run} を使います。

\begin{wslbox}[スクリプトの直接実行]
\begin{lstlisting}[style=bash]
$ cd ~/NMR-lab-kadai/2026_B4_assignment/No.1_FID
$ uv run python makesignal.py
\end{lstlisting}
\end{wslbox}

\section{MATLABとのデータ互換}

このプロジェクトでは \cmd{scipy.io} を通じて
MATLAB の \cmd{.mat} ファイルを読み書きできます。

\begin{lstlisting}[style=python, caption={.mat ファイルの読み込み例}]
import scipy.io

# MATLAB の .mat ファイルを読む
data = scipy.io.loadmat("Signal.mat")

# Python の配列を .mat ファイルに保存する
scipy.io.savemat("result.mat", {"signal": signal})
\end{lstlisting}

\section{インストール済みパッケージの確認}

\begin{wslbox}[パッケージ一覧の表示]
\begin{lstlisting}[style=bash]
$ uv pip list
\end{lstlisting}
\end{wslbox}

\chapter{セットアップ手順まとめ（チートシート）}

\section{最初の1回だけ行う作業}

\begin{wslbox}[(1) uv のインストール]
\begin{lstlisting}[style=bash]
$ curl -LsSf https://astral.sh/uv/install.sh | sh
\end{lstlisting}
\end{wslbox}

ターミナルを閉じて開き直す（または \cmd{source \textasciitilde/.bashrc}）。

\begin{wslbox}[(2) 課題リポジトリの取得]
\begin{lstlisting}[style=bash]
$ cd ~
$ git clone \
  https://github.com/shiru2/NMR-lab-kadai.git
\end{lstlisting}
\end{wslbox}

\begin{wslbox}[(3) 仮想環境の作成・パッケージのインストール]
\begin{lstlisting}[style=bash]
$ cd \
  ~/NMR-lab-kadai/2026_B4_assignment
$ uv sync
\end{lstlisting}
\end{wslbox}

\section{毎回の作業}

\textbf{VS Code（推奨）}：VS Code で
\cmd{2026\_B4\_assignment} フォルダを開くと、
カーネルが自動で \cmd{.venv} に設定されます。

\begin{wslbox}[VS Code でフォルダを開く（ターミナルから）]
\begin{lstlisting}[style=bash]
$ code \
  ~/NMR-lab-kadai/2026_B4_assignment
\end{lstlisting}
\end{wslbox}

\textbf{JupyterLab（任意）}：

\begin{wslbox}[JupyterLab を使う場合]
\begin{lstlisting}[style=bash]
$ cd \
  ~/NMR-lab-kadai/2026_B4_assignment
$ uv run jupyter lab
\end{lstlisting}
\end{wslbox}

ブラウザで \cmd{http://localhost:8888} を開いて作業する。
終了するときはターミナルで \key{Ctrl} + \key{C}。

\section{よく使うコマンドまとめ}

\begin{center}
\begin{tabularx}{\linewidth}{lX}
  \toprule
  コマンド & 内容 \\
  \midrule
  \cmd{uv --version} & uv のバージョン確認 \\
  \cmd{uv sync} & 仮想環境を作成・更新する \\
  \cmd{uv run python スクリプト名} & Python スクリプトを実行する \\
  \cmd{uv run jupyter lab} & JupyterLab を起動する（任意） \\
  \cmd{source .venv/bin/activate} & 仮想環境を手動で有効化する \\
  \cmd{uv pip list} & インストール済みパッケージを確認する \\
  \cmd{uv add パッケージ名} & 新しいパッケージを追加する \\
  \cmd{git pull} & 課題ファイルを最新版に更新する \\
  \cmd{uv self update} & uv 自体をアップデートする \\
  \bottomrule
\end{tabularx}
\end{center}

\section{トラブル時の確認フロー}

課題ノートブックが動かないときは以下の順番で確認してください。

\begin{enumerate}
  \item \cmd{uv sync} を実行してパッケージが揃っているか確認する
  \item VS Code で \cmd{2026\_B4\_assignment} フォルダが
    開かれているか確認する（サブフォルダだけでは不十分）
  \item カーネルエラーが出たら \cmd{rm -rf .venv \&\& uv sync} を試す
  \item それでも解決しない場合は第~\ref{chap:errors}章を参照する
\end{enumerate}

\vfill

\begin{center}
\small
参考：uv 公式ドキュメント \url{https://docs.astral.sh/uv/}
\end{center}

\chapter{研究での活用：uv init / activate / VS Code}

B4課題では既存の \cmd{pyproject.toml} に対して \cmd{uv sync} するだけでしたが、
自分の研究プログラムを書くときは\textbf{ゼロからプロジェクトを作る}必要があります。
この章では研究用途の uv の使い方と、スクリプトを実行する3つの方法、
VS Code との連携を説明します。

% =================================================
\section{uv init：新しいプロジェクトを作る}
% =================================================

自分で研究プログラムを書くときは \cmd{uv init} から始めます。

\begin{wslbox}[新規プロジェクトの作成]
\begin{lstlisting}[style=bash]
$ cd ~
$ uv init myresearch
$ cd myresearch
$ ls
hello.py  pyproject.toml  .python-version  README.md
\end{lstlisting}
\end{wslbox}

\cmd{uv init} は以下のファイルを自動生成します。

\begin{description}
  \item[\cmd{pyproject.toml}]
    プロジェクト設定ファイル（空の依存関係リストで生成される）
  \item[\cmd{.python-version}]
    使用する Python バージョンを固定するファイル
  \item[\cmd{hello.py}]
    サンプルスクリプト（削除して構わない）
\end{description}

\subsection{パッケージを追加する}

\begin{wslbox}[必要なパッケージを追加]
\begin{lstlisting}[style=bash]
$ uv add numpy matplotlib scipy
Resolved 10 packages in 0.3ms
Installed 10 packages in 2.1s
 + numpy==1.26.4
 + matplotlib==3.9.4
 + scipy==1.13.1
\end{lstlisting}
\end{wslbox}

\cmd{uv add} を実行すると：
\begin{itemize}
  \item \cmd{.venv/} が自動作成される
  \item \cmd{pyproject.toml} の \cmd{dependencies} に追記される
  \item \cmd{uv.lock} が自動生成・更新される
\end{itemize}

\subsection{パッケージを削除する}

\begin{wslbox}[パッケージの削除]
\begin{lstlisting}[style=bash]
$ uv remove scipy
\end{lstlisting}
\end{wslbox}

% =================================================
\section{Python スクリプトを実行する 3 つの方法}
% =================================================

仮想環境内の Python でスクリプトを実行する方法は 3 つあります。

\begin{center}
\begin{tabularx}{\linewidth}{l>{\raggedright\arraybackslash}X>{\raggedright\arraybackslash}X}
  \toprule
  方法 & コマンド & 特徴 \\
  \midrule
  \textbf{方法(1)} & \cmd{uv run python script.py} &
    activate 不要・最も推奨 \\
  \textbf{方法(2)} & \cmd{source .venv/bin/activate}
    してから \cmd{python script.py} &
    従来の venv と同じ操作 \\
  \textbf{方法(3)} & \cmd{.venv/bin/python script.py} &
    パスを直接指定・activate 不要 \\
  \bottomrule
\end{tabularx}
\end{center}

\subsection{方法(1) uv run（推奨）}

\begin{wslbox}[方法(1)：uv run]
\begin{lstlisting}[style=bash]
$ uv run python script.py
\end{lstlisting}
\end{wslbox}

\begin{itemize}
  \item activate の必要がない
  \item プロジェクトフォルダ（\cmd{pyproject.toml} がある場所）
    であれば自動的に \cmd{.venv/} を使う
  \item \cmd{uv sync} が済んでいれば問題なく動く
  \item このマニュアルで \textbf{最も推奨する方法}
\end{itemize}

\subsection{方法(2) activate してから実行}

\begin{wslbox}[方法(2)：activate + python]
\begin{lstlisting}[style=bash]
$ source .venv/bin/activate
(.venv) $ python script.py
(.venv) $ python another.py  # activate は一度だけでよい
(.venv) $ deactivate         # 終わったら解除
\end{lstlisting}
\end{wslbox}

\begin{itemize}
  \item \cmd{source .venv/bin/activate} で仮想環境を「有効化」する
  \item 有効化するとプロンプトが \cmd{(.venv) \$} に変わる
  \item 以降は普通に \cmd{python} と打つだけで仮想環境内の Python が使われる
  \item \cmd{deactivate} で元の状態に戻る
  \item 従来の venv 操作に慣れている人はこちらでも問題ない
\end{itemize}

\begin{notebox}[activate のスコープはターミナル1つ]
\cmd{source .venv/bin/activate} の効果は\textbf{そのターミナルウィンドウだけ}です。
新しいターミナルを開いたら、再度 activate する必要があります。
\cmd{uv run} はこの問題がないため、ターミナルをまたぐ作業では便利です。
\end{notebox}

\subsection{方法(3) .venv/bin/python を直接指定}

\begin{wslbox}[方法(3)：パスを直接指定]
\begin{lstlisting}[style=bash]
$ .venv/bin/python script.py
\end{lstlisting}
\end{wslbox}

\begin{itemize}
  \item \cmd{.venv/} 内の Python 実行ファイルを直接パス指定して実行する
  \item activate も \cmd{uv run} も不要
  \item シェルスクリプトや Makefile から呼び出すときに使いやすい
  \item VS Code の Python インタープリタ設定でもこのパスを指定する
\end{itemize}

\begin{wslbox}[どの Python が使われているか確認する]
\begin{lstlisting}[style=bash]
$ which python          # システム Python のパス
/usr/bin/python3

$ uv run which python   # 仮想環境内 Python のパス
/home/username/myresearch/.venv/bin/python

$ .venv/bin/python --version  # 直接確認
Python 3.11.x
\end{lstlisting}
\end{wslbox}

% =================================================
\section{VS Code との連携}
% =================================================

VS Code でプロジェクトを開くと、右下に Python のバージョンが表示されます。
ここで\textbf{正しい仮想環境の Python インタープリタを選ぶ}ことが重要です。

\subsection{インタープリタの選択}

\begin{enumerate}
  \item VS Code でプロジェクトフォルダを開く
    （\cmd{File} $\rightarrow$ \cmd{Open Folder}）
  \item \key{Ctrl} + \key{Shift} + \key{P} でコマンドパレットを開く
  \item \texttt{Python: Select Interpreter} と入力して選択
  \item 一覧に \cmd{.venv/bin/python} が表示されたら選択する
\end{enumerate}

\begin{notebox}[.venv が一覧に出ない場合]
プロジェクトフォルダが VS Code で開かれていないか、
\cmd{uv sync} がまだ完了していない可能性があります。
先に \cmd{uv sync} を実行してから VS Code を開き直してください。
\end{notebox}

\subsection{VS Code の統合ターミナルで実行する}

VS Code 内のターミナル（\cmd{Ctrl} + \cmd{\`{}}）は
インタープリタを設定済みであれば activate 状態で開きます。

\begin{wslbox}[VS Code の統合ターミナルでの実行]
\begin{lstlisting}[style=bash]
(.venv) $ python script.py    # 仮想環境が自動で有効
\end{lstlisting}
\end{wslbox}

\subsection{Python ファイルを「実行」ボタンで動かす}

インタープリタを設定した後、.py ファイルを開いて
右上の \key{$\triangleright$} ボタン（Run ボタン）または \key{F5} で実行できます。
このとき仮想環境が自動的に使われます。

\subsection{Jupyter Notebook を VS Code で開く}

VS Code の Jupyter 拡張機能を使うと、
ブラウザを使わずに直接 \cmd{.ipynb} ファイルを開けます。

\begin{enumerate}
  \item 拡張機能マーケットプレイスで \texttt{Jupyter} をインストール
  \item \cmd{.ipynb} ファイルを VS Code で開く
  \item 右上のカーネル選択で \cmd{.venv/bin/python} を選ぶ
\end{enumerate}

\begin{tipbox}[JupyterLab vs VS Code のどちらを使う？]
\begin{tabularx}{\linewidth}{lX}
  \toprule
  ツール & 向いている場面 \\
  \midrule
  JupyterLab（ブラウザ） &
    グラフをインタラクティブに確認したいとき、
    ノートブック主体の作業 \\
  VS Code &
    .py スクリプトとノートブックを行き来したいとき、
    コード補完・デバッグを活用したいとき \\
  \bottomrule
\end{tabularx}
どちらを使っても仮想環境さえ正しく設定されていれば同じ結果が出ます。
\end{tipbox}

% =================================================
\section{研究プロジェクトの典型的な流れ}
% =================================================

自分の研究プログラムを始めるときの典型的な手順をまとめます。

\begin{wslbox}[研究プロジェクトの開始から実行まで]
\begin{lstlisting}[style=bash]
# 1. プロジェクトを作成
$ uv init myresearch
$ cd myresearch

# 2. 必要なパッケージを追加
$ uv add numpy matplotlib scipy

# 3a. uv run で実行（推奨）
$ uv run python analysis.py

# 3b. activate して実行（従来スタイル）
$ source .venv/bin/activate
(.venv) $ python analysis.py
(.venv) $ deactivate

# 3c. パスを直接指定して実行
$ .venv/bin/python analysis.py
\end{lstlisting}
\end{wslbox}

\begin{notebox}[どれを使うか迷ったら]
\begin{itemize}
  \item 毎回コマンドで実行する $\rightarrow$ \textbf{方法(1)} \cmd{uv run}
  \item VS Code でコーディングしながら実行する $\rightarrow$ \textbf{VS Code} + \textbf{方法(2)}（自動 activate）
  \item シェルスクリプトや自動化に組み込む $\rightarrow$ \textbf{方法(3)} パス直接指定
\end{itemize}
\end{notebox}

% =================================================
\section{よく使う uv コマンド一覧}
% =================================================

\subsection{Python バージョン管理}

\begin{wslbox}[Python バージョンの確認・固定]
\begin{lstlisting}[style=bash]
$ uv python list          # 利用可能な Python の一覧
$ uv python list --only-installed  # インストール済みのみ
$ uv python pin 3.11      # .python-version に 3.11 を固定
\end{lstlisting}
\end{wslbox}

\cmd{uv python pin} を実行すると \cmd{.python-version} ファイルに
バージョンが書き込まれ、以降 \cmd{uv sync} や \cmd{uv run} は
そのバージョンを優先して使います。
共同作業するときや、Python バージョンを明示したいときに使います。

\subsection{依存関係の確認}

\begin{wslbox}[依存関係ツリーの表示]
\begin{lstlisting}[style=bash]
$ uv tree
myresearch v0.1.0
├── matplotlib v3.9.4
│   ├── numpy v1.26.4
│   └── ...
└── scipy v1.13.1
    └── numpy v1.26.4 (*)
\end{lstlisting}
\end{wslbox}

\cmd{uv tree} でどのパッケージが何に依存しているかを確認できます。
\cmd{(*)} は別の場所で既に表示済みであることを示します。

\subsection{requirements.txt へのエクスポート}

pip を使っている人に環境を共有したり、
サーバーにデプロイするときは \cmd{requirements.txt} 形式に変換できます。

\begin{wslbox}[requirements.txt の生成]
\begin{lstlisting}[style=bash]
$ uv export -o requirements.txt
$ uv export --no-hashes -o requirements.txt
\end{lstlisting}
\end{wslbox}

\subsection{ロックファイルだけを更新する}

パッケージをインストールせずに \cmd{uv.lock} だけを
最新の状態に更新したいときは \cmd{uv lock} を使います。

\begin{wslbox}[uv.lock の更新のみ]
\begin{lstlisting}[style=bash]
$ uv lock          # uv.lock を更新（インストールはしない）
$ uv lock --upgrade  # 全パッケージを最新版に更新
\end{lstlisting}
\end{wslbox}

\subsection{uvx：ツールを一時的に実行する}

\cmd{uvx} はインストールせずにツールを一時実行できます（\cmd{pipx run} 相当）。
コードフォーマッタやリンタを手軽に試すときに便利です。

\begin{wslbox}[uvx によるツールの一時実行]
\begin{lstlisting}[style=bash]
$ uvx ruff check .       # コードチェック（ruff）
$ uvx black script.py    # コードフォーマット（black）
$ uvx mypy script.py     # 型チェック（mypy）
\end{lstlisting}
\end{wslbox}

常用するツールはプロジェクトに追加する方法もあります：

\begin{wslbox}[開発ツールをプロジェクトに追加]
\begin{lstlisting}[style=bash]
$ uv add --dev ruff      # 開発依存として追加
$ uv run ruff check .
\end{lstlisting}
\end{wslbox}

\subsection{キャッシュの管理}

uv はダウンロードしたパッケージをキャッシュに保存するため、
2 回目以降のインストールが高速です。
ディスク容量が逼迫したときはキャッシュを削除できます。

\begin{wslbox}[キャッシュの確認と削除]
\begin{lstlisting}[style=bash]
$ uv cache dir    # キャッシュの場所を確認
$ uv cache clean  # キャッシュを全削除
\end{lstlisting}
\end{wslbox}


% ===== 付録 =====
\appendix
\chapter{参考リンク集}

\section{uv 公式ドキュメント}

\url{https://docs.astral.sh/uv/}

英語のみですが、コマンドの使い方や設定の詳細はここが最も正確です。
このマニュアルで扱った内容に対応するページを以下にまとめます。

\section{本マニュアルとの対応}

\begin{center}
\begin{tabularx}{\linewidth}{lX}
  \toprule
  このマニュアルの章 & 対応する公式ドキュメント \\
  \midrule
  第4章 uv のインストール &
    \url{https://docs.astral.sh/uv/getting-started/installation/} \\[0.5em]
  第3章 uv の仕組み（入門） &
    \url{https://docs.astral.sh/uv/getting-started/first-steps/} \\[0.5em]
  第3章 pyproject.toml / uv sync / uv.lock &
    \url{https://docs.astral.sh/uv/concepts/projects/} \\[0.5em]
  第12章 uv init・プロジェクト管理 &
    \url{https://docs.astral.sh/uv/guides/projects/} \\[0.5em]
  第12章 uv run でスクリプト実行 &
    \url{https://docs.astral.sh/uv/guides/scripts/} \\[0.5em]
  第12章 Python バージョン管理 &
    \url{https://docs.astral.sh/uv/concepts/python-versions/} \\[0.5em]
  第8章 JupyterLab との連携 &
    \url{https://docs.astral.sh/uv/guides/integration/jupyter/} \\
  \bottomrule
\end{tabularx}
\end{center}

\begin{notebox}[VS Code との連携について]
uv の公式ドキュメントには VS Code 専用のページはありません。
VS Code 側の設定（Python インタープリタの選択など）は
VS Code 公式ドキュメントを参照してください。\\[0.5em]
\url{https://code.visualstudio.com/docs/python/environments}
\end{notebox}


\end{document}
